\section{Shaping-based low-thrust interplanetary rendezvous}
\label{sec:problem}

\subsection{Low-thrust interplanetary rendezvous}
\label{sec:continuous}

This work focuses on the fixed-time rendezvous problem with
low thrust. The transfer time and number of revolutions are fixed, and the goal
is to minimize an objective function which depends on the control. Since the
spacecraft mass is not included as a state variable, the control
$\mathbf{u}$ represents thrust acceleration rather than thrust force. The
trajectory is parameterized in cylindrical coordinates
$\mathbf{q}=[\rho,\,\theta,\,z]^T$, where $\rho$ is the distance of the spacecraft projection
onto the ecliptic plane from the Sun, $\theta$ is the ecliptic longitude, and
$z$ is the displacement normal to the ecliptic. The heliocentric distance is $r=\sqrt{\rho^2+z^2}$. Using the cylindrical velocity $\mathbf{v}=[v_\rho,\,v_\theta,\,v_z]^T$, the
low-thrust rendezvous problem is

\begin{mini}
	{\mathbf{u}(\cdot)}
	{J(\mathbf{u}(\cdot))}
	{\label{prob:main}}{}
	\addConstraint{\dot{\rho}=v_\rho,\qquad\quad\>\>\>
		\dot{v}_\rho=\frac{v_\theta^2}{\rho}-\frac{\mu\rho}{r^3}+u_\rho}
	\addConstraint{\dot{\theta}=\frac{v_\theta}{\rho},\qquad\quad\>\>
		\dot{v}_\theta=-\frac{v_\rho v_\theta}{\rho}+u_\theta}
	\addConstraint{\dot{z}=v_z,\qquad\quad\>\>\>
		\dot{v}_z=-\frac{\mu z}{r^3}+u_z}
	\addConstraint{\mathbf{q}(t_0)=\mathbf{q}_0,\qquad
		\mathbf{v}(t_0)=\mathbf{v}_0}
	\addConstraint{\mathbf{q}(t_f)=\mathbf{q}_f,\qquad
		\mathbf{v}(t_f)=\mathbf{v}_f}
	\addConstraint{\lVert\mathbf{u}(t)\rVert\leq u_{\max},
		\qquad t\in[t_0,t_f].}
\end{mini}

The boundary states are obtained from the ephemerides of the departure and
target bodies at $t_0$ and $t_f$, respectively. The transfer time
$T=t_f-t_0$ and the number of revolutions $N\in\mathbb{N}$ are fixed. The
unwrapped terminal longitude is selected consistently with $N$, i.e., it
differs from the physical terminal longitude by the required multiple of
$2\pi$.

Two objectives are considered. In the minimum-$\Delta V$ problem,
$J=\Delta V=\int_{t_0}^{t_f}\lVert\mathbf{u}(t)\rVert\,\mathrm{d}t$ for a
prescribed acceleration cap $u_{\max}$; solving it for a sequence of caps
traces the $u_{\max}$--$\Delta V$ Pareto front, and when no cap is prescribed
the path constraint is omitted and the peak acceleration is evaluated a
posteriori. In the minimum-peak problem, $u_{\max}$ is itself a decision
variable and $J=u_{\max}$, so the path constraint is active by construction.
In both cases the path constraint must hold at every instant of the transfer,
which is the difficulty addressed in this work.

\subsection{Polynomial shaping as a finite-dimensional problem}
\label{sec:shaping}

Shaping methods parameterize the spacecraft trajectory with a set of basis
functions. Some examples are Fourier series, B\'ezier curves or B-splines. By
parameterizing the trajectory, the dynamics and the boundary conditions can be
satisfied by design. In this work, polynomial shaping (e.g. B\'ezier curves,
B-splines, etc.) is analyzed. Let us parameterize the cylindrical coordinates as
\begin{equation}
\mathbf{q}(\tau)=\sum^{n_c}_{i=1}\mathbf{c}_i B_i(\tau),\qquad\tau\in[0,1],
\label{eq:shaping}
\end{equation}
where $B_i$ are polynomial basis functions, $\mathbf{c}_i$ are control points,
and $\tau$ is the curve parameter. Throughout this section $\tau=(t-t_0)/T$ is
the normalized time; Section~\ref{sec:sundman} replaces it by a normalized
Sundman arc, and everything below carries over with the corresponding
change of the endpoint rates. With that parameterization, the derivatives of the
cylindrical coordinates are
\begin{equation}
\mathbf{q}'(\tau)=\sum^{n_c}_{i=1}\mathbf{c}_i B_i'(\tau),\quad\quad
\mathbf{q}''(\tau)=\sum^{n_c}_{i=1}\mathbf{c}_i B_i''(\tau),
\end{equation}
where primes denote derivatives with respect to $\tau$, and the physical-time
derivatives follow from the chain rule,
$\dot{\mathbf q}=\mathbf q'/T$ and $\ddot{\mathbf q}=\mathbf q''/T^2$. Then,
the required control acceleration $\mathbf{u}$ can be obtained from the inverse dynamics of
Problem~\eqref{prob:main} as
\begin{equation}
	\mathbf{u}
	=
	\begin{bmatrix}
		u_\rho\\
		u_\theta\\
		u_z
	\end{bmatrix}
	=
	\dfrac{1}{r^3}\begin{bmatrix}
		r^3(\ddot{\rho}-\rho\dot{\theta}^{2})+\mu\rho\\[1mm]
		r^3(2\dot{\rho}\dot{\theta}+\rho\ddot{\theta})\\[1mm]
		r^3\ddot{z}+\mu z
	\end{bmatrix}=\dfrac{1}{W}
	\begin{bmatrix}
	P_\rho\\[1mm]
	P_\theta\\[1mm]
	P_z
	\end{bmatrix},\qquad W=r^3,
	\label{eq:flatness2}
\end{equation}
so that the control is a rational function of the shaped coordinates and
their first two derivatives. This is the flatness property of the low-thrust cylindrical
dynamics. No integration is required and the control is uniquely determined by
the trajectory shape.

In this work the basis functions $B_i$ are those of a clamped B-spline of
degree $n$. On each knot span the shaped coordinates are polynomials of
degree $n$ and can be converted exactly into Bernstein form using a B-spline
extraction matrix. This local Bernstein representation is subsequently used
to construct polynomial products, integrals, and convex-hull bounds.

Substituting Eq.~\eqref{eq:shaping} into Problem~\eqref{prob:main} gives the
finite-dimensional problem
\begin{mini}
	{\mathbf{C}}
	{J(\mathbf{C})}
	{\label{prob:shaped}}{}
	\addConstraint{\mathbf{q}(0;\mathbf C)=\mathbf{q}_0,\qquad
		\mathbf{q}'(0;\mathbf C)=\mathbf{q}'_0}
	\addConstraint{\mathbf{q}(1;\mathbf C)=\mathbf{q}_f,\qquad
		\mathbf{q}'(1;\mathbf C)=\mathbf{q}'_f}
	\addConstraint{\lVert\mathbf{u}(\tau;\mathbf C)\rVert\leq u_{\max},
		\qquad\tau\in[0,1],}
\end{mini}
where $\mathbf{C}=[\mathbf{c}^T_1,\hdots,\mathbf{c}^T_{n_c}]^T$ stacks all the
control points, which are the decision variables of the problem. The dynamics
are satisfied identically through
Eq.~\eqref{eq:flatness2}, and $\mathbf q'_0$, $\mathbf q'_f$ are the
cylindrical rates obtained from $\mathbf v_0$ and $\mathbf v_f$ and expressed
in the curve parameter. Two ingredients remain to be made algebraic: the
rendezvous conditions, which are linear in the control points and are
eliminated in Section~\ref{sec:rendezvous}, and the thrust constraint, which
is semi-infinite and is the subject of Section~\ref{sec:thrust} and of the
remainder of the paper.

\subsubsection{Rendezvous conditions}
\label{sec:rendezvous}

For a clamped B-spline of degree $n$, only the first basis function is
nonzero at $\tau=0$ and only the first two have a nonzero derivative there;
the same holds at $\tau=1$ for the last two. The endpoint values and rates
therefore determine the two control points adjacent to each end,
\begin{equation}
	\mathbf{c}_1=\mathbf{q}_0,\qquad
	\mathbf{c}_2=\frac{\mathbf{q}'_0-B_1'(0)\,\mathbf{c}_1}{B_2'(0)},\qquad
	\mathbf{c}_{n_c}=\mathbf{q}_f,\qquad
	\mathbf{c}_{n_c-1}
	=\frac{\mathbf{q}'_f-B_{n_c}'(1)\,\mathbf{c}_{n_c}}{B_{n_c-1}'(1)}.
	\label{eq:clamped-elimination}
\end{equation}
The remaining $n_c-4$ control points of each coordinate are the free
variables of Problem~\eqref{prob:shaped}, and the rendezvous conditions are
satisfied exactly by construction. The number of revolutions enters only
through the unwrapped terminal longitude $\theta_f$. When the shaped
quantities are not the cylindrical coordinates themselves but primitives
related to them by a nonlinear map, as in Section~\ref{sec:enforcement}, the
prescribed endpoint values and rates are first transformed to the primitives
and Eq.~\eqref{eq:clamped-elimination} is applied to those.

\subsubsection{Thrust constraint}
\label{sec:thrust}

Let us focus on the control acceleration constraint by expanding the shaped
control acceleration of Eq.~\eqref{eq:flatness2},
\begin{equation}
\dfrac{1}{W}\sqrt{P^2_\rho+P^2_\theta+P^2_z}\leq u_{\max}.
\end{equation}
To simplify the previous expression let us square both sides of the
constraint and multiply the resulting expression by $W^2$,
\begin{equation}
P^2_\rho+P^2_\theta+P^2_z \leq W^2 u^2_{\max},
\label{eq:squared-thrust-constraint}
\end{equation}
where the main obstruction to applying polynomial positivity techniques is
that in cylindrical coordinates $W=r^3=(\rho^2+z^2)^{3/2}$ is not a
polynomial. Note that $W$ also appears in $P_\rho$, $P_\theta$ and $P_z$.
Enforcing the inequality at a finite set of nodes is the usual remedy, but it
leaves the constraint unverified between nodes. Section~\ref{sec:enforcement}
removes the obstruction and enforces the bound on the whole trajectory.

\section{Analytical enforcement of the thrust constraint}
\label{sec:enforcement}

\subsection{Levi-Civita map of the meridional plane}
\label{sec:pythagorean}

The obstruction identified in Section~\ref{sec:thrust} is the square root in
$r=\sqrt{\rho^2+z^2}$ which arises when using cylindrical coordinates. For a
fixed $\theta$, the coordinates $(\rho,z)$ describe the spacecraft position
in the meridional plane. To circumvent the obstruction we use the Levi-Civita
map of the meridional plane: writing $\rho+\mathrm{i}z$ as the square of a
complex primitive $w=p+\mathrm{i}q$, the distance becomes
$r=|w^2|=|w|^2$, which is free of square roots. In components, the
Levi-Civita map of the $\rho z$ plane is
\begin{equation}
\begin{aligned}
\rho &= p^2-q^2,\\
z &= 2pq,
\end{aligned}
\label{eq:levi-civita-map}
\end{equation}
and the heliocentric distance is
\begin{equation}
r = \sqrt{\rho^2+z^2}=\sqrt{(p^2+q^2)^2}=p^2+q^2.
\label{eq:levi-civita-radius}
\end{equation}
Then, we switch from the cylindrical coordinates $(\rho,\>\theta,\> z)$ to the
new set of primitives $(p,\>q,\>\theta)$. The primitives $p$ and $q$ encode
the position within the meridional plane. The map is two-to-one, since
$(p,q)$ and $(-p,-q)$ give the same $(\rho,z)$. The ambiguity is removed by a
single sign choice: the cylindrical description requires $\rho>0$, hence
$p^2>q^2$, so $p$ never vanishes along the trajectory and keeps its sign.
Imposing $p>0$ at the departure point therefore fixes the branch on the
whole arc, and the bound $\rho\geq\rho_{\min}>0$ of the optimization problem
guarantees it during the solution. The longitude $\theta$ is retained as a
third primitive, which keeps the number of revolutions an endpoint value as
in Section~\ref{sec:rendezvous}.

If $p$ and $q$ are polynomials of degree $n$ in the curve parameter, it
follows from Eq.~\eqref{eq:levi-civita-radius} that $r$ is a polynomial of
degree $2n$ and $W=r^3$ is a polynomial of degree $6n$.

\subsection{Physical-time thrust certificate}
\label{sec:time-certificate}

With the Levi-Civita map of the meridional plane, the thrust constraint
becomes polynomial whenever the primitives are shaped with polynomial basis
functions. The time derivatives of $\rho$ and $z$ are
\begin{equation}
	\begin{alignedat}{2}
		\dot \rho
		&=2(p\dot p-q\dot q),
		&\qquad
		\dot z
		&=2(\dot p\,q+p\dot q),
		\\
		\ddot \rho
		&=2\left(\dot p^{2}+p\ddot p
		-\dot q^{2}-q\ddot q\right),
		&\qquad
		\ddot z
		&=2\left(\ddot p\,q+2\dot p\dot q+p\ddot q\right).
	\end{alignedat}
	\label{eq:levi-civita-time-derivatives}
\end{equation}
Let us define the kinematic acceleration terms
\begin{equation}
	\begin{aligned}
		\mathcal A_\rho
		&=\ddot{\rho}-\rho\dot{\theta}^2=2\left(\dot p^{2}+p\ddot p
		-\dot q^{2}-q\ddot q\right)-(p^2-q^2)\dot{\theta}^2,\\
		\mathcal A_\theta
		&=2\dot{\rho}\dot{\theta}+\rho\ddot{\theta}=4(p\dot p-q\dot q)\dot{\theta}+(p^2-q^2)\ddot{\theta},\\
		\mathcal A_z
		&=\ddot{z}=2\left(\ddot p\,q+2\dot p\dot q+p\ddot q\right),
	\end{aligned}
	\label{eq:time-kinematic-terms}
\end{equation}
and the cubic heliocentric distance $W=r^3$ as
\begin{equation}
W = (p^2+q^2)^3=p^6+3p^4q^2+3p^2q^4+q^6.
\label{eq:levi-civita-weight}
\end{equation}
Then, the thrust constraint of Eq.~\eqref{eq:squared-thrust-constraint} can
be written in the meridional Levi-Civita primitives and the heliocentric
longitude as
\begin{equation}
	\underbrace{
		\left(W\mathcal A_\rho+\mu (p^2-q^2)\right)^2
	}_{P_\rho^{2}}
	+
	\underbrace{
		\left(W\mathcal A_\theta\right)^2
	}_{P_\theta^{2}}
	+
	\underbrace{
		\left(W\mathcal A_z+2\mu pq\right)^2
	}_{P_z^{2}}
	\leq
	W^2\,u_{\max}^2.
	\label{eq:pythagorean-time-certificate}
\end{equation}
Time derivatives are taken with respect to the normalized time $\tau$, which
only introduces powers of the fixed transfer time $T$ and does not change the
polynomial degrees. If $p$, $q$, and $\theta$ are degree-$n$ polynomials on
each span, the degrees of the terms in
Eq.~\eqref{eq:pythagorean-time-certificate} are given in
Table~\ref{tab:thrust-certificate-degrees}.

\begin{table}[htbp]
	\centering
	\caption{Polynomial degrees of the squared thrust terms for degree-$n$
		Levi-Civita primitives in physical time.}
	\label{tab:thrust-certificate-degrees}
	\begin{tabular}{lc}
		\toprule
		\textbf{Polynomial term} & \textbf{Degree} \\
		\midrule
		$P_\rho^{2}$   & $20n-4$ \\
		$P_\theta^{2}$ & $18n-4$ \\
		$P_z^{2}$      & $16n-4$ \\
		$W^2=r^6$      & $12n$   \\
		\midrule
		$W^2u_{\max}^2-P_\rho^{2}-P_\theta^{2}-P_z^{2}$
		& $20n-4$ \\
		\bottomrule
	\end{tabular}
\end{table}

Requiring every Bernstein coefficient of the degree-$(20n-4)$ polynomial
$W^2u_{\max}^2-P_\rho^{2}-P_\theta^{2}-P_z^{2}$ to be nonnegative on each span
is a sufficient condition for the thrust constraint and constitutes the
physical-time baseline of this work. Its degree, however, makes the condition
expensive and ill-conditioned already for moderate $n$, which motivates the
two developments that follow: a change of independent variable that lowers
the degree of the numerators below that of $W$
(Section~\ref{sec:sundman}), and a convex-hull condition that never forms the
squared polynomial (Section~\ref{sec:hull}).

\subsection{Sundman transformation}
\label{sec:sundman}

The Sundman transformation reduces the polynomial degree that arises when the
primitives are shaped in physical time $t$. It replaces $t$ with the parameter
$s$ through
\begin{equation}
	\mathrm{d}t=r\,\mathrm{d}s,
	\label{eq:sundman-transformation}
\end{equation}
where $r>0$. Denoting physical-time derivatives with dots and derivatives
with respect to the Sundman parameter $s$ with primes, the chain rule gives
\begin{equation}
	\dot f=\frac{f'}{r},
	\qquad
	\ddot f
	=\frac{1}{r}\left(\frac{f'}{r}\right)'
	=\frac{rf''-r'f'}{r^3}.
	\label{eq:sundman-chain-rule}
\end{equation}
The second derivative carries the same denominator $r^3$ as the gravity
term, so no further multiplication by $W=r^3$ is needed to clear the control.
In cylindrical components, substituting Eq.~\eqref{eq:sundman-chain-rule}
into Eq.~\eqref{eq:flatness2} gives
\begin{equation}
\mathbf{u}=\dfrac{1}{W}\begin{bmatrix}
r\rho''-r'\rho'-r\rho\,(\theta')^2+\mu\rho\\
\rho\,(r\theta''-r'\theta')+2r\rho'\theta'\\
rz''-r'z'+\mu z
\end{bmatrix},
\label{eq:sundman-control}
\end{equation}
with the same $W=r^3$ as before. Each entry equals $r^3$ times the physical
kinematic acceleration of Eq.~\eqref{eq:time-kinematic-terms} plus the
gravity numerator, but the factor $r^3$ is now supplied by the chain rule
instead of being multiplied in. Accordingly, in the Sundman parameter the
kinematic terms are redefined as $r^3$ times the physical ones,
$\mathcal A_\rho=r^3(\ddot\rho-\rho\dot\theta^2)$,
$\mathcal A_\theta=r^3(2\dot\rho\dot\theta+\rho\ddot\theta)$ and
$\mathcal A_z=r^3\ddot z$, which in the Levi-Civita primitives read
\begin{equation}
	\begin{aligned}
		\mathcal A_\rho
		&=r\rho''-r'\rho'-r\rho\,(\theta')^2
		=2r\left(pp''-qq''\right)
		-\rho\left[2\left(p'^2+q'^2\right)+r\,(\theta')^2\right],\\
		\mathcal A_\theta
		&=\rho\,(r\theta''-r'\theta')+2r\rho'\theta'
		= r\rho\,\theta''
		+2\left[pp'\left(p^2+3q^2\right)-qq'\left(3p^2+q^2\right)\right]\theta',\\
		\mathcal A_z
		&=rz''-r'z'
		= 2r\left(qp''+pq''\right)-2z\left(p'^2+q'^2\right),
	\end{aligned}
	\label{eq:sundman-kinematic-pq}
\end{equation}
with $r=p^2+q^2$, $\rho=p^2-q^2$ and $z=2pq$. The thrust constraint of
Eq.~\eqref{eq:squared-thrust-constraint} becomes
\begin{equation}
	\underbrace{
		\left(\mathcal A_\rho+\mu (p^2-q^2)\right)^2
	}_{P_\rho^{2}}
	+
	\underbrace{
		\mathcal A_\theta^{2}
	}_{P_\theta^{2}}
	+
	\underbrace{
		\left(\mathcal A_z+2\mu pq\right)^2
	}_{P_z^{2}}
	\leq
	W^2\,u_{\max}^2.
	\label{eq:sundman-certificate}
\end{equation}
Each numerator contains at most one factor of $r$ together with two
derivatives, so $P_\rho$, $P_\theta$ and $P_z$ have degrees $6n-2$, $5n-2$
and $4n-2$, all below $\deg W=6n$. The squared constraint of
Eq.~\eqref{eq:sundman-certificate} has degree $12n$ instead of the $20n-4$
of Eq.~\eqref{eq:pythagorean-time-certificate}.

The change of independent variable alters the curve parameter of
Section~\ref{sec:shaping}. The primitives are now shaped in the normalized
Sundman arc $\sigma=(s-s_0)/S\in[0,1]$, where the total arc
$S=s_f-s_0$ is unknown a priori and becomes an additional decision variable;
derivatives with respect to $s$ are those with respect to $\sigma$ divided by
$S$. The rendezvous conditions of Section~\ref{sec:rendezvous} are imposed
on the primitives in $\sigma$. To preserve the prescribed physical endpoint
velocities, the primitive endpoint derivatives must satisfy
\begin{equation}
	\frac{\mathrm{d}f}{\mathrm{d}\sigma}=S\,r\,\dot f
	\label{eq:sundman-endpoint-rates}
\end{equation}
for every primitive $f$, where $r$ is known at the endpoints, so the
eliminated control points of Eq.~\eqref{eq:clamped-elimination} are affine
in $S$. The prescribed time of flight is enforced through
\begin{equation}
	T=t_f-t_0
	=\int_{s_0}^{s_f}r\,\mathrm{d}s
	=S\int_{0}^{1}r(\sigma)\,\mathrm{d}\sigma,
	\label{eq:normalized-sundman-time-of-flight}
\end{equation}
whose integrand is polynomial on each span, so the constraint is polynomial
in the control points and linear in $S$.
